\section{The observation quotient and the principal theorem}
\label{sec:v81-principal}

An endpoint law may conceal both its generating action and the identity
of the trajectory which produced it.  A calibrated classifier separates
these two problems only by supplying the component identities in another
experiment.  We instead use two noisy readings of the same retained
trajectory.  Neither reading is calibrated, and neither has a prescribed
interpretation as a reflection word.  Variation of the deadline separates
the hidden components.  A third deadline removes even the scalar
normalization of each endpoint law.

The relevant observation is a matrix, not a collection of separately
normalized label laws.  Let $Q\Subset\R^d$ be a compact rectangle and let
$T_1<T_2<T_3$.  For two finite readout alphabets of sizes $u,v$, write
$M_j(z)\in\R^{u\times v}$ for the joint successful density of endpoint
$z$ and the two readings.  Our model is
\begin{equation}\label{eq:v81-model}
 M_j(z)=U(z)\diag\{a_b(z)(T_j-W_b(z))\}_{b=1}^{B}V(z)^{\mathsf T}.
\end{equation}
Here $a_b>0$, $W_b<T_1$, and the columns of $U,V$ are probability
vectors.  The latent number $B$ is not given; $B\le\min(u,v)$ and both
matrices have full column rank.  Conditional on $(b,z)$, the two readings
are independent.  Their conditional laws may depend arbitrarily and
smoothly on $z$, but are the same at every deadline.  This last
invariance is a hypothesis, not a consequence of using two sensors.

The more coarsened observation is
\begin{equation}\label{eq:v81-projective-data}
                         N_j(z)=c_j(z)M_j(z),\qquad c_j(z)>0,
\end{equation}
where none of the functions $c_j$ is known.  This includes a common
branch-blind detector depending on endpoint and deadline, unknown
exposures, normalized endpoint/readout laws, and their combinations.
It also includes observing only the conditional distribution of the
\emph{pair} of readings at each endpoint: then
$c_j(z)=(\one^{\mathsf T}M_j(z)\one)^{-1}$.  Frequencies within the
joint readout matrix are retained; normalizing its rows separately is
not allowed.

\begin{theorem}[Blind clock-pencil rigidity]\label{thm:v81-main}
Assume the preceding model on a connected gate, with $C^m$ coefficients,
$m\ge0$, and distinct actions $W_b(z)$ at every endpoint.

From two matrices $M_1,M_2$, the observations determine $B$, the absolute
actions, the weights, and both readout channels, up to one simultaneous
permutation.  No true-component audit or list of structural branch
identities is supplied.

If $B\ge2$, three projective matrices $N_1,N_2,N_3$ determine $B$, the
absolute actions, both channels and all relative weights $a_b/\sum_c a_c$.
Only a common positive weight multiplier remains unobservable.  Thus
the actions survive arbitrary positive branch-blind detector changes at
each endpoint and deadline, even when all success masses are discarded.

On compact classes with bounded $C^m$ norms, positive weight and residual
floors, positive clock gaps, full-rank margins and separated actions,
these reconstructions are locally Lipschitz in $C^m$.  In the projective
case impose in addition bounds and positive floors for $c_j$, or the
equivalent observed scale and two-plane margins specified below.  The
inverse loses no spatial derivatives.
\end{theorem}

The proof is given in Sections~\ref{sec:v81-affine} and
\ref{sec:v81-projective}.  The simultaneous permutation can be fixed by
increasing action on each connected gate.  This is an intrinsic ordering
of recovered real functions, not a supplied word or lift label.  It
cannot manufacture the semantic names needed by a marked billiard
experiment.  It is unnecessary when one reconstructs a canonical graph
as an unordered union of generating sheets.

The new observation assumption is substantive.  It replaces one
calibrated readout and a ground-truth audit by two conditionally
independent, deadline-invariant readouts.  It permits endpoint-dependent
channels and learns the number of visible components.  It does not
assert that one may obtain independent information by copying a noisy
label or randomly relabelling a single sensor output.  Conditional
independence given the hidden branch and endpoint has to come from the
acquisition mechanism.  Rank, separation and a finite sensor capacity
are also essential to the uniform theorem.

\subsection{Relation to earlier inverse mechanisms}

Spectral identification of latent variables is classical.  Kruskal-type
identifiability, observable-operator methods, repeated measurements and
nonclassical measurement error provide important precedents
\cite{V81AllmanMatiasRhodes,V81AnandkumarHsuKakade,V81BonhommeJochmansRobin,
V81HuSchennach}.  We do not claim a new general uniqueness theorem for
tensor decompositions.  In the present model the known affine dependence
on deadline identifies an \emph{absolute generating action} rather than
an unnamed eigenvalue.  More importantly, three projective observations
recover this affine dependence without measured exposures or a calibrated
channel.  That step makes a branch-blind physical detector an exact
nuisance symmetry instead of a first-order error term.

The deterministic conversion of action to a canonical relation is
classical generating-function calculus
\cite{MarsdenWest2001,LallWest2006}; the conversion of return action to a
Reeb suspension uses the usual return-time primitive
\cite{Hutchings2016,V81Geiges2008}.  The additional conclusion here is upstream of
those constructions: their local absolute action sheets, including the
sheet count, can be recovered from unpaired, uncalibrated projective
records.  Coverage of the required phase domain is still a hypothesis.

For billiards, endpoint actions on a regular gate determine local
covectors by differentiation.  Consequently they contain local canonical
relation information, not merely lengths of periodic trajectories.
A finite endpoint atlas does not determine an entire global scattering
relation unless a coverage theorem supplies the missing trajectories.
Conversely, a scattering relation recorded in Euclidean boundary
coordinates supplies information absent from abstract scalar labels.
We make no unconditional information ordering between these different
experiments.  Marked spectral rigidity results such as
\cite{FinamoreLeguil,DeSimoiKaloshinLeguil2023} and the planar lens theorem
\cite{NoakesStoyanov2020} start from other data and impose other geometric
hypotheses.  The marked/reference-atlas reconstruction retained later in
this article is not relabelled as a canonical unmarked billiard theorem.

\section{A blind inverse for two affine matrix observations}
\label{sec:v81-affine}

Put $\Delta=T_2-T_1$ and
\begin{equation}\label{eq:v81-slope}
 D=\frac{M_2-M_1}{\Delta}=U\diag(a_b)V^{\mathsf T}.
\end{equation}
Since $U,V$ have full column rank and the weights are positive,
$\operatorname{rank}D=B$.  This determines the latent dimension without
knowing either channel.  Fix index sets $I,J$ of $B$ rows and columns
for which $D_{IJ}$ is invertible.  Define the reduced observed operator
\begin{equation}\label{eq:v81-K}
                     K=(M_1)_{IJ}(D_{IJ})^{-1}.
\end{equation}
The invertibility of $D_{IJ}=U_I\diag(a_b)V_J^{\mathsf T}$ implies
that both square factors $U_I,V_J$ are invertible.  Therefore
\begin{equation}\label{eq:v81-diagonal}
                K=U_I\diag(T_1-W_b)U_I^{-1}.
\end{equation}
In particular its eigenvalues are positive and determine the absolute
actions, with multiplicities, by $W_b=T_1-\lambda_b(K)$.
This identity holds even when some actions coincide.  Distinctness is
used for separating the channels and for the stated smooth inverse.

\begin{lemma}[Observable rank-one components]\label{lem:v81-components}
For distinct eigenvalues $\lambda_b=T_1-W_b$, let
\begin{equation}\label{eq:v81-projector}
 P_b=\prod_{c\ne b}\frac{K-\lambda_c\mathrm{Id}}{\lambda_b-\lambda_c},
 \qquad E_b=D_{:J}(D_{IJ})^{-1}P_bD_{I:}.
\end{equation}
Then $E_b=a_bU_{:b}V_{:b}^{\mathsf T}$, independently of the chosen
invertible minor.  In particular
\begin{equation}\label{eq:v81-channels}
 a_b=\one^{\mathsf T}E_b\one,\qquad
 U_{:b}=E_b\one/a_b,\qquad V_{:b}=E_b^{\mathsf T}\one/a_b.
\end{equation}
\end{lemma}
\begin{proof}
Equation~\eqref{eq:v81-diagonal} gives
$P_b=U_I e_be_b^{\mathsf T}U_I^{-1}$.  Also
$D_{:J}(D_{IJ})^{-1}=UU_I^{-1}$ and
$D_{I:}=U_I\diag(a_b)V^{\mathsf T}$.  Substitution proves the
rank-one formula.  Summing its entries and its rows or columns proves
\eqref{eq:v81-channels}, because both channel columns have total mass
one.  The resulting expression depends only on the true component and
not on $I,J$.
\end{proof}

This proves the exact two-clock part of Theorem~\ref{thm:v81-main}.
It also proves uniqueness against alternatives with a different unknown
number of components, provided the alternative satisfies the same full
column-rank hypothesis.  Splitting a component into invisible duplicate
columns is outside that hypothesis and cannot be excluded from data.

\begin{proposition}[Uniform finite-order inverse]\label{prop:v81-stability}
Fix the dimensions, $m$, clock values, a bound $L$ on the $C^m$ norms
of the coefficients and observations, and positive constants
$a_*,r_*,s_*,\gamma$.  Assume
\[
 a_b\ge a_*,\quad T_1-W_b\ge r_*,\quad
 \sigma_B(D)\ge s_*,\quad |W_b-W_c|\ge\gamma\quad(b\ne c).
\]
On a finite cover by minor charts with
$\sigma_{\min}(D_{IJ})\ge\mu_*>0$, formulas
\eqref{eq:v81-K}--\eqref{eq:v81-channels} admit local extensions to
nearby real matrix data.  After fixing the action ordering, their
$C^m$ errors are bounded by $C$ times the $C^m$ observation error.
Here $C$ depends only on the displayed constants, the finite cover,
its fixed gluing functions and $\Delta^{-1}$.  No unknown higher
spatial derivative is used in this assertion.
\end{proposition}
\begin{proof}
Matrix inversion is smooth on the set of matrices with smallest
singular value at least $\mu_*/2$.  The product and reciprocal
inequalities in $C^m$ bound the reduced operator and its perturbation.
For its characteristic polynomial $p$, the derivative at a simple root
has magnitude
$|p'(\lambda_b)|=\prod_{c\ne b}|\lambda_b-\lambda_c|
\ge\gamma^{B-1}$.  The finite-dimensional implicit function theorem
makes each root a smooth real function of the real matrix entries in
a neighborhood of the exact operator.  Its derivatives of every fixed
order are bounded on the compact set under consideration.  Realness
also follows from conjugation and uniqueness of the root in a small
disk about each distinct real root.  Composition with $C^m$ matrix
functions, and repeated differentiation through order $m$, gives the
claimed root estimate.  The projector formula and positive denominators
in~\eqref{eq:v81-channels} give the other estimates.

Compactness and rank provide finitely many minor charts.  One can take
fixed smooth gluing functions subordinate to a slightly smaller cover.
For exact data the local inverses agree by
Lemma~\ref{lem:v81-components}.  For perturbed nonmodel data the glued
outputs need not be an exact latent model, but have the same error bound.
Sorting the separated real actions consistently on the connected gate
prevents a spatially varying permutation.  This construction supplies an
inverse on an open neighborhood, rather than only on the model image.
\end{proof}

\begin{remark}[What can be counted stably]\label{rem:v81-count}
Exact rank determines $B$, not the number of arbitrary hidden causes in
a larger rank-deficient representation.  If a matrix perturbation has
operator norm at most $e<s_*/4$, singular-value thresholding at $s_*/2$
recovers the rank of the reference model.  This follows from the
singular-value perturbation bound.  An extra component whose entire
observed contribution has norm at most $e$ cannot in general be counted
reliably.  The theorem learns visible components above a signal margin;
it does not assert exact open-world discovery at zero signal.
\end{remark}

\section{Three projective observations and unknown detectors}
\label{sec:v81-projective}

The matrices in~\eqref{eq:v81-projective-data} need not have comparable
masses, so their difference is not the slope $D$.  Nevertheless the
known deadlines identify their correct relative scales.  Set
\[
 A=U\diag(a_b)V^{\mathsf T},\qquad
 B_0=U\diag(a_bW_b)V^{\mathsf T};\qquad M_j=T_jA-B_0.
\]
The matrices $A,B_0$ are linearly independent whenever not all $W_b$
are equal.  Indeed $B_0=tA$ would, after multiplication by left and
right inverses of $U,V^{\mathsf T}$, imply $W_b=t$ for every $b$.
Thus any two $M_j$, and hence $N_1,N_2$, are linearly independent.

\begin{lemma}[Clock deprojectivization]\label{lem:v81-projective}
Suppose $N_j=c_j(T_jA-B_0)$, where $c_j>0$ and $A,B_0$ are linearly
independent.  There are uniquely determined scalars $\alpha,\beta$ with
$N_3=\alpha N_1+\beta N_2$.  Put
\begin{equation}\label{eq:v81-rescale}
 \alpha_0=\frac{T_2-T_3}{\Delta},\quad
 \beta_0=\frac{T_3-T_1}{\Delta},\qquad
 \widetilde M_1=\frac{\alpha}{\alpha_0}N_1,\quad
 \widetilde M_2=\frac{\beta}{\beta_0}N_2.
\end{equation}
Then $\widetilde M_j=c_3M_j$ for $j=1,2$.  In particular their
matrix-pencil inverse gives the original absolute actions and channels.
\end{lemma}
\begin{proof}
Affine interpolation in the known variable $T$ gives
$M_3=\alpha_0M_1+\beta_0M_2$.  Hence
$\alpha=\alpha_0c_3/c_1$ and $\beta=\beta_0c_3/c_2$.
Uniqueness follows from linear independence.  Both scale factors in
\eqref{eq:v81-rescale} are positive.  Their substitution proves the
claim.  The common positive multiplier $c_3$ cancels from
\eqref{eq:v81-K}; it multiplies each $E_b$ by $c_3$ and hence cancels
from the recovered channels and relative weights as well.
\end{proof}

To make the computation explicit, vectorize the matrices and form
$H=(\operatorname{vec}N_1\ \operatorname{vec}N_2)$.  Then
\begin{equation}\label{eq:v81-alpha-beta}
 \binom\alpha\beta=(H^{\mathsf T}H)^{-1}
                             H^{\mathsf T}\operatorname{vec}N_3.
\end{equation}
For exact data the projection residual vanishes.  The same expression
is defined on a neighborhood of exact data, even when the residual is
nonzero.  If $\sigma_{\min}(H)\ge\chi_*>0$, the observed norms are
bounded, and the two rescaling factors have a positive floor, it is a
uniformly smooth map there.  Product and reciprocal estimates in $C^m$
followed by Proposition~\ref{prop:v81-stability} prove the projective
stability assertion.  These are the observed scale and two-plane margins
referred to in Theorem~\ref{thm:v81-main}.

The dimension can be recovered before rescaling: since $T_1-W_b>0$,
$\operatorname{rank}N_1=B$.  Equivalently it is the rank of the rescaled
slope.  The equality of these ranks is a model consistency check.  At
least two distinct actions are needed pointwise in the projective
argument.  A single constant action with arbitrary unknown exposures
is not identified by any number of such scalar projective observations.

\begin{corollary}[Exact immunity to a common detector]
\label{cor:v81-detector}
Suppose a retention decision made before the two readouts has probability
$e_j(z)>0$ conditional on endpoint and deadline, independently of the
latent component.  The functions $e_j$ are unknown, may differ between
clocks, and need not be close to one another.  Under the projective
hypotheses of Theorem~\ref{thm:v81-main}, they cause no ambiguity in
any absolute action or recovered channel.  The same is true if the
failure counts or the total exposure at each deadline are discarded.
\end{corollary}
\begin{proof}
The detector multiplies every entry of the matrix by the same scalar
$e_j(z)$.  Unknown exposure and normalization add further positive scalar
factors.  They are all absorbed in $c_j$ in
\eqref{eq:v81-projective-data}.  Apply
Lemma~\ref{lem:v81-projective} and the blind affine inverse.
\end{proof}

Thus the common-detector model has an exact nonlinear quotient, not an
$O(\epsilon)$ efficiency floor.  This does not contradict a lower bound
for independently branch-dependent retention: that larger apparatus
class does not act by scalar multiplication of the observed matrices.
The distinction will be made physical in
Section~\ref{sec:v81-physical}.

\begin{proposition}[Two projective clocks are insufficient]
\label{prop:v81-two-projective}
Even with two independent full-rank readouts, distinct smooth actions
and strict margins, two projective clocks do not generally determine
absolute actions.  There is a nontrivial one-parameter ambiguity in the
same class.
\end{proposition}
\begin{proof}
Let $r_b=(T_2-W_b)/(T_1-W_b)>1$.  For $q>0$ sufficiently close to one,
put
\[
 W_b^{(q)}=\frac{T_2-qr_bT_1}{1-qr_b},\qquad
 a_b^{(q)}=a_b\frac{T_1-W_b}{T_1-W_b^{(q)}}.
\]
Then $a_b^{(q)}(T_1-W_b^{(q)})=a_b(T_1-W_b)$ and
$a_b^{(q)}(T_2-W_b^{(q)})=q a_b(T_2-W_b)$.  Keep $U,V$ unchanged,
replace $c_2$ by $c_2/q$ and keep $c_1$ fixed.  Both observations are
unchanged.  Choose the initial weights small enough that the total raw success
masses are strictly below one.  Compactness then preserves the raw
probability constraints, positivity, distinctness and every strict
margin for $q$ close to one.  The fractional map is not the
identity unless $q=1$, so the action tuple changes.  This obstruction
is in the unrestricted generating-action observation class; no claim
that every transformed action is a billiard action is required.
\end{proof}

\begin{proposition}[The joint second readout carries information]
\label{prop:v81-one-view}
After discarding the second reading, even known raw masses at arbitrarily
many deadlines need not identify the action tuple in an uncalibrated
one-readout model.  Full column rank of the unknown one-readout channel
does not repair this.
\end{proposition}
\begin{proof}
Take a strictly positive full-column-rank stochastic matrix $U$, positive
$a$, and distinct $W$.  The observed vector is $g_j=U(T_ja-aW)$, where
$aW$ denotes the componentwise product.  Let $R$ be an invertible matrix
close to the identity with $\one^{\mathsf T}R=\one^{\mathsf T}$, not a
permutation.  Put $U'=UR$, $a'=R^{-1}a$, and
$a'W'=R^{-1}(aW)$.  For a sufficiently small nontrivial perturbation,
$U'$ remains strictly positive and stochastic, $a'>0$, and the actions
remain distinct and admissible.  All vectors $g_j$ agree for every
clock.  One may choose $R$ so that at least one action changes, for
example by mixing two unequal action components.  Two independently
noisy readings are therefore not merely two names for one observation.
\end{proof}
